%chktex-file 1 chktex-file 2 chktex-file 3
%chktex-file 8 chktex-file 9 chktex-file 10
%chktex-file 13 %chktex-file 17 %chktex-file 23
%chktex-file 45 %chktex-file 25 %chktex-file 35
%chktex-file 40

\documentclass[oneside,12pt,a4paper]{report}

\usepackage{template}

\begin{document}

    \begin{center}
        \fontsize{17.28}{17.28}
            \textbf{SME0340 -- Equações Diferenciais Ordinárias}

            \textbf{Terceira Lista de Exercícios}
    \end{center}

%%%%%%%%%%%%%%%%%%%%%%%%%%%%%%%%%%%%%%%%%%%%%%%%%%%%%%%%%
% QUESTÕES
%%%%%%%%%%%%%%%%%%%%%%%%%%%%%%%%%%%%%%%%%%%%%%%%%%%%%%%%%

\exercicio[Zill, Seção 4.1, Exercícios 7, 8] 
\begin{itensquestao}
    \item Dado que $x(t) = c_1 \cos \omega t + c_2 \sin \omega t$ é a solução geral
    de $x'' + \omega^2 x = 0$ no intervalo $(-\infty, \infty)$, mostre que uma solução
    que satisfaça as condições iniciais $x(0) = x_0$ e $x'(0) = x_1$ é dada por
    \[x(t) = x_0 \cos \omega t + \frac{x_1}{\omega} \sin \omega t.\]

    \item Use a solução geral de $x'' + \omega^2 x = 0$ dada no Problema 7 para
    mostrar que uma solução satisfazendo as condições $x(t_0) = x_0$ e
    $x'(t_0) = x_1$ é a solução dada no Problema 7, deslocada por $t_0$ unidades:
    \[x(t) = x_0 \cos \omega(t - t_0) + \frac{x_1}{\omega} \sin\, \omega(t - t_0).\]
\end{itensquestao}

\exercicio[Zill, Seção 4.1, Exercícios 16, 17, 20, 22] Determine se o conjunto de funções dado é
linearmente independente no intervalo $(-\infty, \infty)$.
\begin{itensquestao}[1]
    \item $f_1(x) = 0, \quad f_2(x) = x, \quad f_3(x) = e^x$
    \item $f_1(x) = 5, \quad f_2(x) = \cos^2 x, \quad f_3(x) = \sin^2 x$
    \item $f_1(x) = 2 + x, \quad f_2(x) = 2 + |x|$
    \item $f_1(x) = e^x, \quad f_2(x) = e^{-x}, \quad f_3(x) = \sinh x$
\end{itensquestao}

\exercicio[Zill, Seção 4.1, Exercício 35]
\begin{itensquestao}
\item Verifique que $y_{p_1} = 3e^{2x}$ e $y_{p_2} = x^2 + 3x$ são,
respectivamente, soluções particulares de
\begin{align*}
    y'' - 6y' + 5y &= -9e^{2x} \\
    \text{e} \quad y'' - 6y' + 5y &= 5x^2 + 3x - 16.
\end{align*}

\item Use o item (a) para encontrar soluções particulares de
\begin{align*}
    y'' - 6y' + 5y &= 5x^2 + 3x - 16 - 9e^{2x} \\
    \text{e} \quad y'' - 6y' + 5y &= -10x^2 - 6x + 32 + e^{2x}.
\end{align*}
\end{itensquestao}

\exercicio[Zill, Seção 4.2, Exercícios 17, 19, 20] A função indicada $y_1(x)$ é uma solução da equação homogênea associada. Use o método de redução de ordem para encontrar uma segunda solução $y_2(x)$ da equação homogênea e uma solução particular da equação não-homogênea dada.
\begin{itensquestao}
    \item $y''-4y=2,\quad y_1=e^{-2x}$
    \item $y''-3y'+2y=5e^{3x},\quad y_1=e^{x}$
    \item $y''-4y'+3y=x,\quad y_1=e^{x}$
\end{itensquestao}

\exercicio[Zill, Seção 4.3, Exercícios 29, 32, 35, 36] Resolva o problema de valor inicial dado.
\begin{itensquestao}
    \item $y'' + 16y = 0, \quad y(0) = 2, \quad y'(0) = -2$
    \item $4y'' - 4y' - 3y = 0, \quad y(0) = 1, \quad y'(0) = 5$
    \item $y''' + 12y'' + 36y' = 0, \quad y(0) = 0, \quad y'(0) = 1, \quad y''(0) = -7$
    \item $y''' + 2y'' - 5y' - 6y = 0, \quad y(0) = y'(0) = 0, \quad y''(0) = 1$
\end{itensquestao}

\exercicio[Zill, Seção 4.4, Exercícios 14, 17] Use coeficientes a determinar para resolver a equação diferencial dada.
\begin{itensquestao}[2]
    \item $y'' - 4y = (x^2 - 3) \sin 2x$
    \item $y'' - 2y' + 5y = e^x \cos 2x$
\end{itensquestao}

\exercicio[Zill, Seção 4.4, Exercício 41] Resolva o problema de valor inicial no qual a função entrada $g(x)$ é descontínua. [\textit{Sugestão}: Resolva o problema em dois intervalos e então ache uma solução de tal forma que $y$ e $y'$ sejam contínuas em $x = \pi/2$.]
$$y'' + 4y = g(x), \quad y(0) = 1, \quad y'(0) = 2,~\text{onde}~g(x) = \begin{cases} \sin x, & 0 \leq x \leq \frac{\pi}{2} \\ 0, & x > \frac{\pi}{2} \end{cases}$$

\exercicio[Zill, Seção 4.5, Exercícios 11, 13] Verifique se o operador diferencial dado anula as funções indicadas.
\begin{itensquestao}
    \item $D^4; \quad y = 10x^3 - 2x$
    \item $(D - 2)(D + 5); \quad y = e^{2x} + 3e^{-5x}$
\end{itensquestao}

\exercicio[Zill, Seção 4.5, Exercícios 45, 49, 55] Resolva a equação diferencial dada usando coeficientes a determinar.
\begin{itensquestao}
    \item $y''-2y'-3y=4e^x-9$
    \item $y''+6y'+9y=-xe^{4x}$
    \item $y''+25y=20\sin(5x)$
\end{itensquestao}

\exercicio[Zill, Seção 4.6, Exercícios 9, 12, 15, 17 e 21] Resolva cada equação diferencial por variação dos parâmetros, respeitando as condições iniciais quando houver.
\begin{itensquestao}[2]
    \item $y''-4y=\dfrac{e^{2x}}{x}$
    \item $y''-2y'+y=\dfrac{e^{x}}{1+x^2}$
    \item $y''+2y'+y=e^{-t}\ln(t)$
    \item $3y''-6y'+6y=e^x\sec(x)$
    \item $y''+2y'-8y=2e^{-2x}-e^{-x},\\ y(0)=1,~y'(0)=0$
\end{itensquestao}

\exercicio[Zill, Seção 4.6, Exercício 23] As funções indicadas são soluções linearmente independentes da equação diferencial homogênea associada em $(0,\infty)$. Ache a solução geral da equação não homogênea dada:
\[x^2y''+xy'+\left(x^2-\frac{1}{4}\right)y=x^\frac{3}{2};\quad y_1=x^{-\frac{1}{2}}\cos(x),~y_2=x^{-\frac{1}{2}}\sin(x)\]

\exercicio[Zill, Seção 4.6, Exercícios 25, 27] Resolva as equações diferenciais de terceira ordem por variação dos parâmetros.
\begin{itensquestao}[2]
    \item $y'''+y'=\tan(x)$
    \item $y'''-2y''-y'+2y=e^{4x}$
\end{itensquestao}

\exercicio[] Resolva a equação diferencial
\[y''+7y=\sum_{k=3}^{n}e^{3kt}\]

\exercicio[] Um circuito em série tem um gerador $E(t)=500~\mathrm{volts}$, um capacitor de $\frac{1}{108}\times10^{-3}~\mathrm{farad}$, um resistor de $300~\mathrm{ohms}$ e um indutor de $0.2~\mathrm{henry}$. Se a carga inicial no capacitor é de $\frac{217}{216}~\mathrm{coulomb}$ e não há corrente inicial, encontre a carga $Q(t)$ no capacitor e a corrente no circuito $I(t)$ no instante $t$.
    

%%%%%%%%%%%%%%%%%%%%%%%%%%%%%%%%%%%%%%%%%%%%%%%%%%%%%%%%%
% RESOLUÇÕES
%%%%%%%%%%%%%%%%%%%%%%%%%%%%%%%%%%%%%%%%%%%%%%%%%%%%%%%%%

\begin{resolucao}
    \itemresolucao Primeiramente, derivemos a solução geral fornecida:
        \[x'(t)=-\omega c_1\sin(\omega t)+\omega c_2\cos(\omega t)\]

    Agora, devemos substituir as condições iniciais fornecidas:
    \[\begin{cases}
        x(0)=x_0\Rightarrow \boxed{x_0=c_1}\\
        x'(0)=x_1\Rightarrow x_1=\omega c_2\Rightarrow \boxed{c_2=\dfrac{x_1}{\omega}}
    \end{cases}\]

    Substituindo as constantes na solução geral, temos
    \[\boxed{x(t)=x_0\cos(\omega t)+\frac{x_1}{\omega}\sin(\omega t)}\]

    \itemresolucao O procedimento é análogo ao do item anterior. Aproveitando a primeira derivada obtida, substituímos as condições iniciais fornecidas:
    \[\begin{cases}
        x(t_0)=x_0\Rightarrow x_0=c_1\cos(\omega t)+c_2\sin(\omega t)\\
        x'(t_0)=x_1\Rightarrow x_1=-\omega c_1\sin(\omega t)+\omega c_2\cos(\omega t)
    \end{cases}\]

    Ou, na forma matricial, 
    \[\begin{Bmatrix}
        x_0\\
        x_1
    \end{Bmatrix}
    =
    \begin{pmatrix}
        \cos(\omega t_0) & \sin(\omega t_0)\\
        -\omega\sin(\omega t_0) & \omega\cos(\omega t_0)
    \end{pmatrix}
    \begin{Bmatrix}
        c_1\\
        c_2
    \end{Bmatrix}\]

    Resolvendo o sistema,
    \[\begin{Bmatrix}
        c_1\\
        c_2
    \end{Bmatrix}
    =
    \underbrace{\frac{1}{\omega\cos^2(\omega t)+\omega\sin^2(\omega t)}}_{\omega}
    \begin{pmatrix}
        \omega\cos(\omega t_0) & -\sin(\omega t_0)\\
        \omega\sin(\omega t_0) & \cos(\omega t_0)
    \end{pmatrix}
    \begin{Bmatrix}
        x_0\\
        x_1
    \end{Bmatrix}\]
    \[\begin{cases}
        c_1=x_0\cos(\omega t_0)-\dfrac{x_1}{\omega}\sin(\omega t_0)\\
        c_2=x_0\sin(\omega t_0)+\dfrac{x_1}{\omega}\cos(\omega t_0)
    \end{cases}\]

    Substituindo as constantes encontradas na solução geral fornecida, temos
    \[x(t)=\left(x_0\cos(\omega t_0)-\dfrac{x_1}{\omega}\sin(\omega t_0)\right)\cos(\omega t)+\left(x_0\sin(\omega t_0)+\dfrac{x_1}{\omega}\cos(\omega t_0)\right)\sin(\omega t)\]
    \[x(t)=x_0(\cos(\omega t_0)\cos(\omega t)+\sin(\omega t_0)\sin(\omega t))-\dfrac{x_1}{\omega}(\sin(\omega t_0)\cos(\omega t)-\cos(\omega t_0)\sin(\omega t))\]
    \[\boxed{x(t)=x_0\cos(\omega(t-t_0))-\dfrac{x_1}{\omega}\sin(\omega(t-t_0))}\]
\end{resolucao}

\begin{resolucao}
    Usaremos duas definições para resolver os itens dessa questão:
    \begin{enumerate}
        \item \textbf{Wronskiano:} Suponha que cada uma das funções $f_1(x), f_2(x), \ldots, f_n(x)$ tenha
        pelo menos $n - 1$ derivadas. O determinante

        $$W(f_1, f_2, \ldots, f_n) =
        \begin{vmatrix}
        f_1            & f_2            & \cdots & f_n            \\
        f_1'           & f_2'           & \cdots & f_n'           \\
        \vdots         & \vdots         & \cdots & \vdots         \\
        f_1^{(n-1)}    & f_2^{(n-1)}    & \cdots & f_n^{(n-1)}
        \end{vmatrix},$$

        onde as linhas denotam derivadas, é chamado de \textbf{wronskiano das funções}.

        \item \textbf{Critério de independência linear:} Sejam $y_1, y_2, \ldots, y_n$ $n$ soluções da equação diferencial linear
        homogênea de ordem $n$ em um intervalo $I$. Então, o conjunto de soluções
        será \textbf{linearmente independente} em $I$ se e somente se
        $W(y_1, y_2, \ldots, y_n) \neq 0$ para todo $x$ no intervalo.
    \end{enumerate}

    \itemresolucao 
    \[W(x)=\begin{array}{|ccc|}
        0 & x & e^x\\
        0 & 1 & e^x\\
        0 & 0 & e^x
    \end{array}=0\quad\forall~x\in(-\infty, \infty)\]
    Portanto, as soluções são linearmente dependentes.

    \itemresolucao 
    \[W(x)=\begin{array}{|ccc|}
        5 & \cos^2(x) & \sin^2(x)\\
        0 & -\sin(2x) & \sin(2x)\\
        0 & -2\cos(2x) & 2\cos(2x)
    \end{array}=0\quad\forall~x\in(-\infty, \infty)\]
    Portanto, as soluções são linearmente dependentes.

    \itemresolucao A função $f_2(x)$ não é diferenciável em $x=0$. Assim, o Wronskiano não é definido para todo $x\in(-\infty, \infty)$ e, consequentemente, as soluções são linearmente dependentes.

    \itemresolucao 
    \[W(x)=\begin{array}{|ccc|}
        e^x & e^{-x} & \sinh(x)\\
        e^x & -e^{-x} & \cosh(x)\\
        e^x & e^{-x} & \sinh(x)
    \end{array}\]
    \[W(x)=e^x\cdot
    \begin{array}{|cc|}
        -e^{-x} & \cosh(x)\\
        e^{-x} & \sinh(x)
    \end{array}
    -
    e^x\cdot
    \begin{array}{|cc|}
        e^{-x} & \sinh(x)\\
        e^{-x} & \sinh(x)
    \end{array}
    +
    e^x\cdot
    \begin{array}{|cc|}
        e^{-x} & \sinh(x)\\
        e^{-x} & \cosh(x)
    \end{array}
    \]
    \[W(x)=0\quad\forall~x\in(-\infty, \infty)\]
    Portanto, as soluções são linearmente dependentes.
\end{resolucao}

\begin{resolucao}
    \itemresolucao Para verificar que $y_{p_1}$ é solução particular da equação dada, devemos substituir $y_{p_1}$ na equação e verificar se a igualdade é satisfeita. Derivando $y_{p_1}$, temos
    \[\begin{cases}
        y_{p_1}'=6e^{2x}\\
        y_{p_1}''=12e^{2x}
    \end{cases}\]

    Dessa forma,
    \[12e^{2x}-36e^{2x}+15e^{2x}=-9e^{2x}\]
    \[\boxed{-9e^{2x}=-9e^{2x}}\]

    Como a equação foi reduzida à identidade, confirmamos que $y_{p_1}$ é solução particular da primeira equação dada. De maneira análoga, derivando $y_{p_2}$, temos
    \[\begin{cases}
        y_{p_2}'=2x+3\\
        y_{p_2}''=2
    \end{cases}\]

    Assim,
    \[2-12x-18+5x^2+15x=5x^2+3x-16\]
    \[\boxed{5x^2+3x-16=5x^2+3x-16}\]

    Como a equação foi reduzida à identidade, confirmamos que $y_{p_2}$ é solução particular da segunda equação dada.

    \itemresolucao Dado um operador diferencial linear $L$ e um termo não homogêneo $g(x)$, se $L[y_{p_1}]=g_1(x)$ e $L[y_{p_2}]=g_2(x)$, então $L[c_1y_{p_1}+c_2y_{p_2}]=c_1L[y_{p_1}]+c_2L[y_{p_2}]=c_1g_1(x)+c_2g_2(x)$. Com base no item anterior, definimos $g_1(x)=-9e^{2x}$ e $g_2(x)=5x^2+3x-16$. Dessa forma, para a primeira equação dada, temos
    \[y''-6y'+5y=1 \cdot g_2(x) + 1 \cdot g_1(x)\Rightarrow y_p(x)=1\cdot y_{p_2}+1\cdot y_{p_1}\]
    \[\boxed{y_p(x)=x^2+3x+3e^{2x}}\]
    Para a segunda equação,
    \[y''-6y'+5y=-2 \cdot g_2(x) + \frac{1}{9} \cdot g_1(x)\Rightarrow y_p(x)=-2\cdot y_{p_2}+\frac{1}{9}\cdot y_{p_1}\]
    \[\boxed{y_p(x)=-2x^2-6x-\frac{1}{3}e^{2x}}\]
\end{resolucao}

\begin{resolucao}
    \itemresolucao Realizando a substituição $y(x)=u(x)y_1(x)=u(x)e^{-2x}$, temos:
    \[\begin{cases}
        y'=u'y_1+uy_1'=-2e^{-2x}u+e^{-2x}u'\\
        y''=y_1''u+2y'u'+y_1u''=4e^{-2x}u-4e^{-2x}u'+e^{-2x}u''
    \end{cases}\]

    Substituindo na equação diferencial fornecida, 
    \[4e^{-2x}u-4e^{-2x}u'+e^{-2x}u''-4e^{-2x}u=2\]
    \[e^{-2x}(u''-4u')=2\]
    \[u''-4u'=2e^{2x}\]

    A partir da substituição $w=u'$, a equação é reduzida para a primeira ordem. Assim,
    \[w'-4w=2e^{2x}\]

    Um fator integrante para essa equação é $\mu(x)=e^{\int -4dx}=e^{-4x}$. Multiplicando a equação por $\mu(x)$, temos
    \[\frac{d}{dt}\left(we^{-4x}\right)=2e^{-2x}\Rightarrow we^{-4x}=-e^{-2x}+c_2\]
    \[w=-e^{2x}+c_2'e^{4x}\]

    Integrando $w$ para obter $u$, temos
    \[u=-\frac{1}{2}e^{2x}+\frac{c_2'}{4}e^{4x}+c_1\]

    Substituindo de volta em $y$,
    \[y(x)=e^{-2x}\left(-\frac{1}{2}e^{2x}+\frac{c_2'}{4}e^{4x}+c_1\right)\]
    \[y(x)=-\frac{1}{2}+c_2e^{2x}+c_1e^{-2x}\]

    Onde $c_2=c_2'/4$. Assim, a segunda solução da equação homogênea associada é $\boxed{y_2(x)=e^{2x}}$ e uma solução particular da equação não homogênea é $\boxed{y_p(x)=-\frac{1}{2}}$.

    \itemresolucao Realizando a substituição $y(x)=u(x)y_1(x)=u(x)e^{x}$, temos:
    \[\begin{cases}
        y'=u'y_1+uy_1'=e^{x}u+e^{x}u'\\
        y''=y_1''u+2y'u'+y_1u''=e^{x}u+e^{x}u'+e^{x}u''
    \end{cases}\]

    Substituindo na equação diferencial fornecida, 
    \[e^{x}u+e^{x}u'+e^{x}u''-3e^{x}u-3e^{x}u'+2e^{x}u=0\]
    \[e^xu''-e^xu'=5e^{3x}\]
    \[u''-u'=5e^{2x}\]

    A partir da substituição $w=u'$, a equação é reduzida para a primeira ordem. Assim,
    \[w'-w=5e^{2x}\]

    Um fator integrante para essa equação é $\mu(x)=e^{\int -1dx}=e^{-x}$. Multiplicando a equação por $\mu(x)$, temos
    \[\frac{d}{dt}\left(we^{-x}\right)=5e^{x}\Rightarrow we^{-x}=5e^{x}+c_2\]
    \[w=5e^{2x}+c_2e^{x}\]

    Integrando $w$ para obter $u$, temos
    \[u=\frac{5}{2}e^{2x}+c_2e^{x}+c_1\]

    Substituindo de volta em $y$,
    \[y(x)=e^{x}\left(\frac{5}{2}e^{2x}+c_2e^{x}+c_1\right)\]
    \[y(x)=\frac{5}{2}e^{3x}+c_2e^{2x}+c_1e^{x}\]

    Assim, a segunda solução da equação homogênea associada é $\boxed{y_2(x)=e^{2x}}$ e uma solução particular da equação não homogênea é $\boxed{y_p(x)=\frac{5}{2}e^{3x}}$.

    \itemresolucao Realizando a substituição $y(x)=u(x)y_1(x)=u(x)e^{x}$, temos:
    \[\begin{cases}
        y'=u'y_1+uy_1'=e^{x}u+e^{x}u'\\
        y''=y_1''u+2y'u'+y_1u''=e^{x}u+2e^{x}u'+e^{x}u''
    \end{cases}\]

    Substituindo na equação diferencial fornecida, 
    \[e^{x}u+2e^{x}u'+e^{x}u''-4e^{x}u-4e^{x}u'+3e^{x}u=0\]
    \[e^xu''-2e^xu'=x\]
    \[u''-2u'=xe^{-x}\]

    A partir da substituição $w=u'$, a equação é reduzida para a primeira ordem. Assim,
    \[w'-2w=xe^{-x}\]

    Um fator integrante para essa equação é $\mu(x)=e^{\int -2dx}=e^{-2x}$. Multiplicando a equação por $\mu(x)$, temos
    \[\frac{d}{dt}\left(we^{-2x}\right)=xe^{-3x}\Rightarrow we^{-2x}=-\frac{1}{3}xe^{-3x}-\frac{1}{9}e^{-3x}+c_2'\]
    \[w=-\frac{1}{3}xe^{-x}-\frac{1}{9}e^{-x}+c_2'e^{2x}\]

    Integrando $w$ para obter $u$, temos
    \[u=\frac{1}{3}xe^{-x}+\frac{4}{9}e^{-x}+\frac{c_2'}{2}e^{2x}+c_1\]

    Substituindo de volta em $y$,
    \[y(x)=e^{x}\left(\frac{1}{3}e^{-x}+\frac{4}{9}e^{-x}+\frac{c_2'}{2}e^{2x}+c_1\right)\]
    \[y(x)=\frac{1}{3}x+\frac{4}{9}+c_2e^{3x}+c_1e^{x}\]

    Onde $c_2=c_2'/2$. Assim, a segunda solução da equação homogênea associada é $\boxed{y_2(x)=e^{3x}}$ e uma solução particular da equação não homogênea é $\boxed{y_p(x)=\frac{1}{3}x+\frac{4}{9}}$.
\end{resolucao}

\begin{resolucao}
    \itemresolucao Para a equação diferencial linear homogênea de segunda ordem $y'' + 16y = 0$, a equação característica associada é dada por:
    \[\lambda^2 + 16 = 0 \Rightarrow \lambda = \pm 4i\]

    Como as raízes são complexas conjugadas da forma $\lambda = \alpha \pm \beta i$, com $\alpha = 0$ e $\beta = 4$, a solução geral é expressa em termos de senos e cossenos:
    \[y(x) = c_1 \cos(4x) + c_2 \sin(4x)\]

    Derivando a solução geral para aplicar as condições iniciais, obtemos:
    \[y'(x) = -4c_1 \sin(4x) + 4c_2 \cos(4x)\]

    Impondo a primeira condição inicial $y(0) = 2$:
    \[y(0) = c_1 \cos(0) + c_2 \sin(0) = 2 \Rightarrow c_1 = 2\]

    Impondo a segunda condição inicial $y'(0) = -2$:
    \[y'(0) = -4(2)\sin(0) + 4c_2 \cos(0) = -2 \Rightarrow 4c_2 = -2 \Rightarrow c_2 = -\frac{1}{2}\]

    Substituindo as constantes encontradas na solução geral, o problema de valor inicial possui a seguinte solução:
    \[\boxed{y(x) = 2 \cos(4x) - \frac{1}{2} \sin(4x)}\]

    \itemresolucao Para a equação diferencial linear homogênea $4y'' - 4y' - 3y = 0$, a equação característica associada é dada por:
    \[4\lambda^2 - 4\lambda - 3 = 0\]

    Resolvendo, encontramos as raízes reais e distintas:
    \[\lambda = \frac{-(-4) \pm \sqrt{(-4)^2 - 4(4)(-3)}}{2(4)} = \frac{4 \pm \sqrt{16 + 48}}{8} = \frac{4 \pm 8}{8}\]
    \[\lambda_1 = \frac{3}{2} \quad \text{e} \quad \lambda_2 = -\frac{1}{2}\]

    Assim, a solução geral da equação diferencial é:
    \[y(x) = c_1 e^{\frac{3}{2}x} + c_2 e^{-\frac{1}{2}x}\]

    Sua primeira derivada é dada por:
    \[y'(x) = \frac{3}{2}c_1 e^{\frac{3}{2}x} - \frac{1}{2}c_2 e^{-\frac{1}{2}x}\]

    Aplicando as condições iniciais fornecidas, montamos o seguinte sistema linear:
    \[\begin{cases}
        y(0) = c_1 + c_2 = 1 \\
        y'(0) = \frac{3}{2}c_1 - \frac{1}{2}c_2 = 5
    \end{cases}\]

    Multiplicando a segunda equação por $2$, obtemos $3c_1 - c_2 = 10$. Somando esta com a primeira equação ($c_1 + c_2 = 1$), temos:
    \[4c_1 = 11 \Rightarrow c_1 = \frac{11}{4}\]

    Substituindo de volta na primeira equação, encontramos $c_2$:
    \[c_2 = 1 - \frac{11}{4} = -\frac{7}{4}\]

    Portanto, a solução do problema de valor inicial é:
    \[\boxed{y(x) = \frac{11}{4}e^{\frac{3}{2}x} - \frac{7}{4}e^{-\frac{1}{2}x}}\]

    \itemresolucao Para a equação diferencial linear homogênea de terceira ordem $y''' + 12y'' + 36y' = 0$, determinamos a equação característica correspondente:
    \[\lambda^3 + 12\lambda^2 + 36\lambda = 0\]

    Colocando $\lambda$ em evidência, temos:
    \[\lambda(\lambda^2 + 12\lambda + 36) = 0 \Rightarrow \lambda(\lambda + 6)^2 = 0\]

    As raízes obtidas são $\lambda_1 = 0$ (simples) e $\lambda_2 = \lambda_3 = -6$ (raiz real de multiplicidade 2). Desse modo, a solução geral assume a forma:
    \[y(x) = c_1 + c_2 e^{-6x} + c_3 x e^{-6x}\]

    Para aplicar as condições iniciais, calculamos a primeira e a segunda derivadas de $y(x)$:
    \[y'(x) = -6c_2 e^{-6x} + c_3 e^{-6x} - 6c_3 x e^{-6x}\]
    \[y''(x) = 36c_2 e^{-6x} - 12c_3 e^{-6x} + 36c_3 x e^{-6x}\]

    Aplicando as condições iniciais no ponto $x = 0$, obtemos o sistema:
    \[\begin{cases}
        y(0) = c_1 + c_2 = 0 \\
        y'(0) = -6c_2 + c_3 = 1 \\
        y''(0) = 36c_2 - 12c_3 = -7
    \end{cases}\]

    Da segunda equação, isolamos $c_3 = 1 + 6c_2$. Substituindo na terceira equação:
    \[36c_2 - 12(1 + 6c_2) = -7 \Rightarrow -36c_2 - 12 = -7 \Rightarrow -36c_2 = 5 \Rightarrow c_2 = -\frac{5}{36}\]

    Consequentemente, determinamos $c_1$ e $c_3$:
    \[c_1 = -c_2 = \frac{5}{36}\]
    \[c_3 = 1 + 6\left(-\frac{5}{36}\right) = 1 - \frac{5}{6} = \frac{1}{6}\]

    Assim, a solução final do problema de valor inicial é:
    \[\boxed{y(x) = \frac{5}{36} - \frac{5}{36}e^{-6x} + \frac{1}{6}xe^{-6x}}\]

    \itemresolucao Para a equação diferencial linear homogênea de terceira ordem $y''' + 2y'' - 5y' - 6y = 0$, escrevemos a equação característica associada:
    \[\lambda^3 + 2\lambda^2 - 5\lambda - 6 = 0\]

    Por inspeção, notamos que $\lambda = -1$ é uma raiz. Dividindo o polinômio por $(\lambda + 1)$ por Briot-Ruffini, obtemos:
    \begin{table}[H]
        \centering
        \begin{tabular}{c|cccc}
            2 & 1 & 2 & -5 & 4\\
            \hline
              & 1 & 4 & 3 & 0
        \end{tabular}
    \end{table}

    Ou seja, $(\lambda+1)(\lambda^2+4\lambda+3)=0=(\lambda+1)(\lambda-2)(\lambda+3)$. As raízes são reais e distintas: $\lambda_1 = -1$, $\lambda_2 = -3$ e $\lambda_3 = 2$. Portanto, a solução geral é dada por:
    \[y(x) = c_1 e^{-x} + c_2 e^{-3x} + c_3 e^{2x}\]

    Calculando as sucessivas derivadas para a aplicação das condições iniciais:
    \[y'(x) = -c_1 e^{-x} - 3c_2 e^{-3x} + 2c_3 e^{2x}\]
    \[y''(x) = c_1 e^{-x} + 9c_2 e^{-3x} + 4c_3 e^{2x}\]

    Impondo as condições iniciais no ponto $x = 0$, obtemos o seguinte sistema linear:
    \[\begin{cases}
        c_1 + c_2 + c_3 = 0 \\
        -c_1 - 3c_2 + 2c_3 = 0 \\
        c_1 + 9c_2 + 4c_3 = 1
    \end{cases}\]

    Somando a primeira e a segunda equações, eliminamos $c_1$:
    \[-2c_2 + 3c_3 = 0 \Rightarrow c_2 = \frac{3}{2}c_3\]

    Subtraindo a primeira equação da terceira equação, também eliminamos $c_1$:
    \[8c_2 + 3c_3 = 1\]

    Substituindo $c_2 = \frac{3}{2}c_3$ na expressão acima:
    \[8\left(\frac{3}{2}c_3\right) + 3c_3 = 1 \Rightarrow 12c_3 + 3c_3 = 1 \Rightarrow 15c_3 = 1 \Rightarrow c_3 = \frac{1}{15}\]

    Encontrando os valores correspondentes para $c_2$ e $c_1$:
    \[c_2 = \frac{3}{2}\left(\frac{1}{15}\right) = \frac{1}{10}\]
    \[c_1 = -c_2 - c_3 = -\frac{1}{10} - \frac{1}{15} = -\frac{5}{30} = -\frac{1}{6}\]

    Substituindo os valores das constantes, a solução do problema de valor inicial é:
    \[\boxed{y(x) = -\frac{1}{6}e^{-x} + \frac{1}{10}e^{-3x} + \frac{1}{15}e^{2x}}\]
\end{resolucao}

\begin{resolucao}
    \itemresolucao Para determinar a solução geral da equação diferencial linear não homogênea, primeiramente resolvemos a equação homogênea associada $y'' - 4y = 0$. A equação característica é dada por:
    \[\lambda^2 - 4 = 0 \Rightarrow \lambda = \pm 2\]

    Como as raízes são reais e distintas, a solução homogênea é:
    \[y_h(x) = c_1 e^{2x} + c_2 e^{-2x}\]

    O termo não homogêneo é $g(x) = (x^2 - 3)\sin(2x)$. Visto que $\pm 2i$ não são raízes da equação característica, propomos uma solução particular da forma:
    \[y_p(x) = (Ax^2 + Bx + C)\cos(2x) + (Dx^2 + Ex + F)\sin(2x)\]

    Calculando a segunda derivada de $y_p(x)$ e substituindo na expressão $y_p'' - 4y_p$, agrupamos os termos em relação a $\cos(2x)$ e $\sin(2x)$:
    \[[ -8Ax^2 + (8D - 4B)x + (2A + 4E - 8C) ]\cos(2x) +\] 
    \[+[ -8Dx^2 - (8A + 4E)x + (2D - 4B - 8F)]\sin(2x) =\]
    \[=(x^2 - 3)\sin(2x)\]

    Igualando os coeficientes correspondentes em ambos os lados da equação, obtemos o seguinte sistema linear:
    \[\begin{cases}
        -8A = 0 \Rightarrow A = 0 \\
        8D - 4B = 0 \Rightarrow B = 2D \\
        2A + 4E - 8C = 0 \\
        -8D = 1 \Rightarrow D = -\frac{1}{8} \\
        -(8A + 4E) = 0 \Rightarrow E = 0 \\
        2D - 4B - 8F = -3
    \end{cases}\]

    A partir de $D = -\frac{1}{8}$, encontramos $B = -\frac{1}{4}$. Substituindo $A=0$ e $E=0$ na terceira equação, obtemos $C = 0$. Por fim, substituindo os valores de $D$ e $B$ na última equação:
    \[2\left(-\frac{1}{8}\right) - 4\left(-\frac{1}{4}\right) - 8F = -3 \Rightarrow -\frac{1}{4} + 1 - 8F = -3 \Rightarrow 8F = \frac{15}{4} \Rightarrow F = \frac{15}{32}\]

    Dessa forma, a solução particular é $\boxed{y_p(x) = -\frac{1}{4}x\cos(2x) + \left(-\frac{1}{8}x^2 + \frac{15}{32}\right)\sin(2x)}$, e a solução geral da equação diferencial é:
    \[\boxed{y(x) = c_1 e^{2x} + c_2 e^{-2x} - \frac{1}{4}x\cos(2x) + \left(-\frac{1}{8}x^2 + \frac{15}{32}\right)\sin(2x)}\]

    \itemresolucao Primeiramente, determinamos a solução da equação homogênea associada $y'' - 2y' + 5y = 0$ por meio de sua equação característica:
    \[\lambda^2 - 2\lambda + 5 = 0 \Rightarrow \lambda = \frac{2 \pm \sqrt{4 - 20}}{2} = 1 \pm 2i\]

    Como as raízes são complexas conjugadas, a solução homogênea correspondente é:
    \[y_h(x) = e^x(c_1 \cos(2x) + c_2 \sin(2x))\]

    O termo forçante é $g(x) = e^x \cos(2x)$. Como o número complexo $1 + 2i$ coincide com uma das raízes da equação característica, ocorre o fenômeno de ressonância. Portanto, devemos multiplicar a forma padrão por $x$, propondo a solução particular:
    \[y_p(x) = x e^x (A \cos(2x) + B \sin(2x))\]

    Para simplificar a substituição, podemos escrever $y_p(x) = e^x z(x)$, onde $z(x) = x(A\cos(2x) + B\sin(2x))$. Inserindo essa forma na equação diferencial $y_p'' - 2y_p' + 5y_p = e^x \cos(2x)$, a expressão reduz-se elegantemente a:
    \[e^x(z'' + 4z) = e^x \cos(2x) \Rightarrow z'' + 4z = \cos(2x)\]

    Calculando as derivadas de $z(x)$:
    \[z'(x) = (A\cos(2x) + B\sin(2x)) + x(-2A\sin(2x) + 2B\cos(2x))\]
    \[z''(x) = 4B\cos(2x) - 4A\sin(2x) - 4x(A\cos(2x) + B\sin(2x))\]

    Substituindo $z''(x)$ em $z'' + 4z = \cos(2x)$, os termos multiplicados por $x$ se cancelam, restando apenas:
    \[4B\cos(2x) - 4A\sin(2x) = \cos(2x)\]

    Comparando os coeficientes, obtemos diretamente:
    \[4B = 1 \Rightarrow B = \frac{1}{4} \quad \text{e} \quad -4A = 0 \Rightarrow A = 0\]

    Desse modo, a solução particular é $\boxed{y_p(x) = \frac{1}{4}xe^x\sin(2x)}$, estabelecendo a solução geral da equação diferencial como:
    \[\boxed{y(x) = e^x(c_1 \cos(2x) + c_2 \sin(2x)) + \frac{1}{4}xe^x\sin(2x)}\]
\end{resolucao}

\begin{resolucao}
    Para resolver o problema de valor inicial com termo forçante descontínuo, dividimos a análise em dois intervalos de continuidade da função entrada $g(x)$.

    \textbf{Intervalo 1:} Para $0 \leq x \leq \frac{\pi}{2}$, a equação diferencial é dada por $y'' + 4y = \sin x$. 
    A equação característica associada à parte homogênea é:
    \[\lambda^2 + 4 = 0 \Rightarrow \lambda = \pm 2i\]
    Fornecendo a solução homogênea $y_h(x) = c_1 \cos(2x) + c_2 \sin(2x)$. Pelo método dos coeficientes a determinar, assumimos uma solução particular da forma $y_p(x) = A \cos x + B \sin x$. Substituindo $y_p(x)$ e sua segunda derivada na EDO, obtemos:
    \[3A \cos x + 3B \sin x = \sin x \Rightarrow A = 0 \quad \text{e} \quad B = \frac{1}{3}\]
    Desse modo, a solução geral no primeiro intervalo é:
    \[y_1(x) = c_1 \cos(2x) + c_2 \sin(2x) + \frac{1}{3} \sin x\]
    Determinamos as constantes $c_1$ e $c_2$ aplicando as condições iniciais em $x = 0$:
    \[y_1(0) = c_1 = 1\]
    \[y_1'(0) = 2c_2 + \frac{1}{3} = 2 \Rightarrow c_2 = \frac{5}{6}\]
    Portanto, a solução para o primeiro intervalo é:
    \[y_1(x) = \cos(2x) + \frac{5}{6} \sin(2x) + \frac{1}{3} \sin x\]
    Para assegurar a continuidade da solução e de sua primeira derivada na fronteira $x = \frac{\pi}{2}$, avaliamos os limites à esquerda:
    \[y_1\left(\frac{\pi}{2}\right) = \cos(\pi) + \frac{5}{6}\sin(\pi) + \frac{1}{3}\sin\left(\frac{\pi}{2}\right) = -\frac{2}{3}\]
    \[y_1'\left(\frac{\pi}{2}\right) = -2\sin(\pi) + \frac{5}{3}\cos(\pi) + \frac{1}{3}\cos\left(\frac{\pi}{2}\right) = -\frac{5}{3}\]

    \textbf{Intervalo 2:} Para $x > \frac{\pi}{2}$, o termo forçante torna-se nulo, resultando na EDO homogênea $y'' + 4y = 0$. A solução geral neste intervalo é:
    \[y_2(x) = c_3 \cos(2x) + c_4 \sin(2x)\]
    Sua derivada é dada por $y_2'(x) = -2c_3 \sin(2x) + 2c_4 \cos(2x)$. Aplicando as condições de colagem em $x = \frac{\pi}{2}$ para garantir que $y$ e $y'$ sejam contínuas:
    \[y_2\left(\frac{\pi}{2}\right) = -c_3 = -\frac{2}{3} \Rightarrow c_3 = \frac{2}{3}\]
    \[y_2'\left(\frac{\pi}{2}\right) = -2c_4 = -\frac{5}{3} \Rightarrow c_4 = \frac{5}{6}\]
    Assim, a solução particular do problema para cada intervalo é escrita como:
    \[\boxed{y(x) = \begin{cases} \cos(2x) + \frac{5}{6} \sin(2x) + \frac{1}{3} \sin x, & 0 \leq x \leq \frac{\pi}{2} \\ \frac{2}{3} \cos(2x) + \frac{5}{6} \sin(2x), & x > \frac{\pi}{2} \end{cases}}\]
\end{resolucao}

\begin{resolucao}
    \itemresolucao O operador diferencial $D^n$ representa a aplicação sucessiva de $n$ derivadas ordinárias em relação à variável independente $x$. Para verificar se o operador $D^4$ anula a função polinomial $y(x) = 10x^3 - 2x$, calculamos suas derivadas sequencialmente:

    A primeira derivada ($D$) resulta em:
    \[D(10x^3 - 2x) = 30x^2 - 2\]

    A segunda derivada ($D^2$) é obtida derivando o resultado anterior:
    \[D^2(10x^3 - 2x) = D(30x^2 - 2) = 60x\]

    A terceira derivada ($D^3$) corresponde a:
    \[D^3(10x^3 - 2x) = D(60x) = 60\]

    Finalmente, a quarta derivada ($D^4$) da função é a derivada de uma constante:
    \[D^4(10x^3 - 2x) = D(60) = 0\]

    Como o resultado final é identicamente nulo, conclui-se que \boxed{\text{o operador } D^4 \text{ anula a função dada}}.

    \itemresolucao Para verificar se o operador composto $(D - 2)(D + 5)$ anula a função $y(x) = e^{2x} + 3e^{-5x}$, podemos aplicar os operadores linearmente e de forma sequencial (da direita para a esquerda).

    Primeiramente, aplicamos o operador interno $(D + 5)$ sobre a função $y(x)$:
    \[(D + 5)(e^{2x} + 3e^{-5x}) = D(e^{2x} + 3e^{-5x}) + 5(e^{2x} + 3e^{-5x})\]

    Efetuando a diferenciação do primeiro termo, obtemos:
    \[D(e^{2x} + 3e^{-5x}) = 2e^{2x} - 15e^{-5x}\]

    Substituindo de volta na expressão do operador interno e agrupando os termos semelhantes:
    \[(2e^{2x} - 15e^{-5x}) + (5e^{2x} + 15e^{-5x}) = 7e^{2x}\]

    Agora, aplicamos o operador restante $(D - 2)$ ao resultado obtido no passo anterior:
    \[(D - 2)(7e^{2x}) = D(7e^{2x}) - 2(7e^{2x})\]
    \[= 14e^{2x} - 14e^{2x} = 0\]

    Como a aplicação sucessiva dos operadores lineares reduziu a expressão a zero, confirma-se que \boxed{\text{o operador } (D - 2)(D + 5) \text{ de fato anula a função dada}}.
\end{resolucao}

\begin{resolucao}

    \itemresolucao Primeiramente, devemos resolver a equação homogênea auxiliar $y'' - 2y' - 3y = 0$. Sua equação característica é
    \[\lambda^2 - 2\lambda - 3 = 0 \Rightarrow (\lambda - 3)(\lambda + 1) = 0 \Rightarrow \lambda_1 = 3,\quad \lambda_2 = -1\]
    Logo, a solução complementar é
    \[y_c(x) = c_1 e^{3x} + c_2 e^{-x}\]

    Em seguida, identificamos o operador anulador de $g(x) = 4e^x - 9$. Como $e^x$ é anulado por $(D-1)$ e a constante $9$ é anulada por $D$, o anulador de $g(x)$ é $L_1 = D(D-1)$. Aplicando $L_1$ a ambos os lados da equação não homogênea, obtemos
    \[D(D-1)(D^2-2D-3)\,y = 0 \Rightarrow D(D-1)(D-3)(D+1)\,y = 0\]

    A solução geral da equação homogênea de ordem superior associada tem raízes $r = 0,\, 1,\, 3,\, -1$, resultando em
    \[y = c_1 e^{3x} + c_2 e^{-x} + c_3 e^{x} + c_4\]

    Desconsiderando os termos que já aparecem em $y_c$, a forma da solução particular é
    \[y_p(x) = Ae^x + B\]

    Calculamos as derivadas necessárias, $y_p' = Ae^x$ e $y_p'' = Ae^x$, e substituímos na equação original:
    \[Ae^x - 2Ae^x - 3(Ae^x + B) = 4e^x - 9\]
    \[-4Ae^x - 3B = 4e^x - 9\]

    Agrupando os coeficientes de cada função, chegamos ao sistema
    \[\begin{cases}
        -4A = 4\\
        -3B = -9
    \end{cases}\Rightarrow A = -1,\quad B = 3\]

    Portanto, $y_p(x) = -e^x + 3$, e a solução geral da equação diferencial é
    \[\boxed{y(x) = c_1 e^{3x} + c_2 e^{-x} - e^x + 3}\]

    \itemresolucao Primeiramente, devemos resolver a equação homogênea auxiliar $y'' + 6y' + 9y = 0$. Sua equação característica é
    \[\lambda^2 + 6\lambda + 9 = 0 \Rightarrow (\lambda + 3)^2 = 0 \Rightarrow \lambda = -3\quad \text{(raiz dupla)}\]
    Como a equação possui raiz repetida, a fim de encontrar soluções linearmente independentes, temos $y_1 = e^{-3x}$ e $y_2 = xe^{-3x}$. Assim, a solução complementar é
    \[y_c(x) = c_1 e^{-3x} + c_2 x e^{-3x}\]

    Em seguida, identificamos o operador anulador de $g(x) = -xe^{4x}$. Como $xe^{4x}$ é um polinômio de grau 1 multiplicado por $e^{4x}$, o operador que anula tal função é $L_1 = (D - 4)^2$. Aplicando $L_1$ a ambos os lados da equação não homogênea, obtemos
    \[(D-4)^2(D+3)^2\,y = 0\]

    As raízes da equação homogênea de ordem superior resultante são $r = 4$ (dupla) e $r = -3$ (dupla), de modo que
    \[y = c_1 e^{-3x} + c_2 x e^{-3x} + c_3 e^{4x} + c_4 x e^{4x}\]

    Desconsiderando os termos que já aparecem em $y_c$, a forma da solução particular é
    \[y_p(x) = Ae^{4x} + Bxe^{4x}\]

    Calculamos as derivadas necessárias:
    \[y_p' = (4A + B + 4Bx)e^{4x}, \qquad y_p'' = (16A + 8B + 16Bx)e^{4x}\]

    Substituindo na equação original e agrupando os termos semelhantes:
    \[(16A + 8B + 16Bx)e^{4x} + 6(4A + B + 4Bx)e^{4x} + 9(A + Bx)e^{4x} = -xe^{4x}\]
    \[(49A + 14B)e^{4x} + 49Bxe^{4x} = -xe^{4x}\]

    Igualando os coeficientes de cada função, chegamos ao sistema
    \[\begin{cases}
        49A + 14B = 0\\
        49B = -1
    \end{cases}\Rightarrow B = -\frac{1}{49},\quad A = \frac{2}{343}\]

    Portanto, $y_p(x) = \dfrac{2}{343}e^{4x} - \dfrac{1}{49}xe^{4x}$, e a solução geral da equação diferencial é
    \[\boxed{y(x) = c_1 e^{-3x} + c_2 x e^{-3x} + \frac{2}{343}e^{4x} - \frac{1}{49}xe^{4x}}\]

    \itemresolucao Primeiramente, devemos resolver a equação homogênea auxiliar $y'' + 25y = 0$. Sua equação característica é
    \[\lambda^2 + 25 = 0 \Rightarrow \lambda = \pm\, 5i\]
    Dessa forma, a solução complementar é
    \[y_c(x) = c_1\cos(5x) + c_2\sin(5x)\]

    Em seguida, identificamos o operador anulador de $g(x) = 20\sin(5x)$. Como $\sin(5x)$ é anulado por $(D^2 + 25)$, o operador anulador é $L_1 = D^2 + 25$. Aplicando $L_1$ a ambos os lados da equação não homogênea, obtemos
    \[(D^2 + 25)^2\,y = 0\]

    As raízes da equação homogênea de ordem superior resultante são $r = \pm\,5i$, cada uma de multiplicidade 2. Assim, a solução geral é
    \[y = c_1\cos(5x) + c_2\sin(5x) + c_3 x\cos(5x) + c_4 x\sin(5x)\]

    Desconsiderando os termos que já aparecem em $y_c$, a forma da solução particular é
    \[y_p(x) = Ax\cos(5x) + Bx\sin(5x)\]

    Note que este é um caso de \textit{ressonância}, visto que as funções $\cos(5x)$ e $\sin(5x)$ já integram $y_c$, exigindo a multiplicação por $x$. Calculamos as derivadas necessárias:
    \[y_p' = (A + 5Bx)\cos(5x) + (B - 5Ax)\sin(5x)\]
    \[y_p'' = (10B - 25Ax)\cos(5x) + (-10A - 25Bx)\sin(5x)\]

    Substituindo na equação original:
    \[y_p'' + 25y_p = 10B\cos(5x) - 10A\sin(5x) = 20\sin(5x)\]

    Igualando os coeficientes de cada função, chegamos ao sistema
    \[\begin{cases}
        10B = 0\\
        -10A = 20
    \end{cases}\Rightarrow B = 0,\quad A = -2\]

    Portanto, $y_p(x) = -2x\cos(5x)$, e a solução geral da equação diferencial é
    \[\boxed{y(x) = c_1\cos(5x) + c_2\sin(5x) - 2x\cos(5x)}\]

\end{resolucao}

\begin{resolucao}
    \itemresolucao Primeiramente, devemos resolver a equação homogênea auxiliar $y''-4y=0$. Sua equação característica é $\lambda^2-4=0$, por meio da qual obtemos $\lambda=\pm2$. Ou seja, as soluções que satisfazem a equação homogênea são $y_1(x) = e^{2x}$ e $y_2(x) = e^{-2x}$. Podemos escrever a solução da equação homogênea como uma combinação linear de $y_1$ e $y_2$, da forma
    \[y_h(x)=c_1e^{2x}+c_2e^{-2x}\]

    Em seguida, prosseguimos com a determinação da solução particular. Em um primeiro momento, calculamos o Wronskiano das soluções:
    \[W(x)=\begin{vmatrix}
        e^{2x} & e^{-2x}\\
        2e^{2x} & -2e^{-2x}
    \end{vmatrix}=-2-2=-4\]

    Além disso,
    \[W_1(x)=\begin{vmatrix}
        0 & e^{-2x}\\
        1 & -2e^{-2x}
    \end{vmatrix}=-e^{-2x}\quad e \quad W_2(x)=\begin{vmatrix}
        e^{2x}  & 0\\
        2e^{2x} & 1
    \end{vmatrix}=e^{2x}\]

    Dessa forma, segue da fórmula de variação de parâmetros que
    \[y_p(x)=\sum_{i=1}^{2}y_i\int\frac{f(x)W_i(x)}{W(x)}dx, \text{ onde } f(x)=\frac{e^{2x}}{x}\]
    \[y_p(x)=e^{2x}\int\frac{e^{2x}(-e^{-2x})}{4x}dx+e^{-2x}\int\frac{e^{2x}e^{2x}}{4x}dx\]
    \[y_p(x)=-\frac{e^{2x}}{4}\int\frac{1}{x}dx+\frac{e^{-2x}}{4}\int\frac{e^{4x}}{x}dx\]

    Enquanto a primeira integral é direta, a segunda é não elementar. Dessa forma, vamos expressá-la por meio da variável muda $t$. Além disso, tomamos um intervalo que não engloba a descontinuidade em $x=0$:
    \[y_p(x)=-\frac{e^{2x}}{4}\ln(x)+\frac{e^{-2x}}{4}\int_{x_0}^{x}\frac{e^{4t}}{t}dt;\quad x_0>0\]

    Temos que $y(x)=y_h(x)+y_p(x)$. Assim,
    \[\boxed{y(x)=c_1e^{2x}+c_2e^{-2x}-\frac{e^{2x}}{4}\ln(x)+\frac{e^{-2x}}{4}\int_{x_0}^{x}\frac{e^{4t}}{t}dt;\quad x_0>0}\]

    \itemresolucao Tomando a equação característica da equação homogênea auxiliar, 
    \[\lambda^2-2\lambda+1=0\Rightarrow(\lambda-1)^2=0\Rightarrow\lambda=1\]
    Como a equação tem raízes repetidas, a fim de encontrar soluções linearmente independentes, temos $y_1=e^x$ e $y_2=xe^x$. Assim, a solução da equação homogênea será uma combinação linear de $y_1$ e $y_2$, da forma
    \[y_h=c_1e^x+c_2xe^x\]
    Calculando o Wronskiano das soluções,
    \[W(x)=\begin{vmatrix}
        e^{x} & xe^{x}\\
        e^{x} & e^{x}(x+1)
    \end{vmatrix}=e^{2x}(1+x)-xe^{2x}=e^{2x}\]

    Além disso,
    \[W_1(x)=\begin{vmatrix}
        0 & xe^{x}\\
        1 & e^{x}(x+1)
    \end{vmatrix}=-xe^{x}\quad e \quad W_2(x)=\begin{vmatrix}
        e^{x}  & 0\\
        e^{x} & 1
    \end{vmatrix}=e^{x}\]

    Segue de fórmula de variação dos parâmetros que
    \[y_p(x)=\sum_{i=1}^{2}y_i\int\frac{f(x)W_i(x)}{W(x)}dx, \text{ onde } f(x)=\frac{e^{x}}{1+x^2}\]
    \[y_p(x)=e^x\int\frac{\frac{e^x}{1+x^2}(-xe^x)}{e^{2x}}dx+xe^x\int\frac{\frac{e^x}{1+x^2}(e^x)}{e^{2x}}dx\]
    \[y_p(x)=-e^x\int\frac{x}{1+x^2}dx+xe^x\int\frac{1}{1+x^2}dx\]
    \[y_p(x)=-\frac{e^x}{2}\ln(1+x^2)+xe^x\arctan(x)\]

    Assim, a solução geral da equação diferencial é dada por $y(x)=y_h(x)+y_p(x)$,
    \[\boxed{y(x)=c_1e^x+c_2xe^x+e^x\left(x\arctan(x)-\frac{\ln(1+x^2)}{2}\right)}\]

    \itemresolucao Resolvendo a equação característica da homogênea auxiliar,
    \[\lambda^2+2\lambda+1=0\Rightarrow(\lambda+1)^2\Rightarrow\lambda=-1\]
    Dessa forma, tomando soluções linearmente independentes, temos $y_1=e^{-t}$ e $y_2=te^{-t}$. Ou seja, a solução da equação homogênea formada pela combinação linear de $y_1$ e $y_2$ é
    \[y_h(t)=c_1e^{-t}+c_2te^{-t}\]

    Partindo para a determinação da solução particular por meio da variação de parâmetros, calculamos o Wronskiano:
    \[W(t)=\begin{vmatrix}
        e^{-t} & te^{-t}\\
        -e^{-t} & e^{-t}(1-t)
    \end{vmatrix}=e^{-2t}(1-t)+te^{-2t}=e^{-2t}\]

    Ademais,
    \[W_1(t)=\begin{vmatrix}
        0 & te^{-t}\\
        1 & e^{-t}(1-t)
    \end{vmatrix}=-te^{-t}\quad e \quad W_2(t)=\begin{vmatrix}
        e^{-t}   & 0\\
        -e^{-t} & 1
    \end{vmatrix}=e^{-t}\]

    Dessa forma, segue que
    \[y_p(t)=e^{-t}\int\frac{e^{-t}\ln(t)(-te^{-t})}{e^{-2t}}dt+te^{-t}\int\frac{e^{-t}\ln(t)e^{-t}}{e^{-2t}}dt\]
    \[y_p(t)=-e^{-t}\int t\ln(t)dt+te^{-t}\int\ln(t)dt\]
    Calcularemos as duas integrais por partes. Para a primeira, tomamos $\ln(t)=u\rightarrow1/t~dt=du$ e $t~dt=dv\rightarrow t^2/2=v$. Já para a segunda, temos $\ln(t)=u\rightarrow1/t~dt=du$ e $dt=dv\rightarrow t=v$. Nesse sentido,
    \[y_p(t)=-e^{-t}\left[\frac{t^2}{2}\ln(t)-\int\frac{t}{2}dt\right]+te^{-t}\left[t\ln(t)-\int dt\right]\]
    \[y_p(t)=-e^{-t}\left[\frac{t^2}{2}\ln(t)-\frac{t^2}{4}\right]+te^{-t}\left[t\ln(t)-t\right]=e^{-t}\left[\frac{t^2}{2}\ln(t)-\frac{3}{4}t^2\right]\]

    Assim, a solução geral da equação é
    \[\boxed{c_1e^{-t}+c_2te^{-t}+e^{-t}\left[\frac{t^2}{2}\ln(t)-\frac{3}{4}t^2\right]}\]

    \itemresolucao Primeiramente, devemos reescrever a equação diferencial na forma padrão:
    \[y''-2y'+2=\frac{e^x\sec(x)}{3}\]
    Resolvendo a equação característica da homogênea auxiliar,
    \[\lambda^2-2\lambda+2=0\Rightarrow\lambda_{1,2}=\frac{2\pm\sqrt{2^2-4\cdot1\cdot2}}{2}\Rightarrow\lambda_{1,2}=1\pm i\]

    Tomando $\lambda=1+i$, temos que as soluções reaisi linearmente independentes da equação são $y_1=e^x\cos(x)$ e $y_2=e^x\sin(x)$. Ou seja, uma combinação linear dessas duas soluções é
    \[y_h(x)=c_1e^x\cos(x)+c_2e^x\sin(x)\]

    Calculemos o Wronskiano das soluções:
    \[W(x)=\begin{vmatrix}
        e^{x}\cos(x) & e^{x}\sin(x)\\
        e^{x}(\cos(x)-\sin(x)) & e^{x}(\sin(x)+\cos(x))
    \end{vmatrix}=e^{2x}\]

    Ademais,
    \[W_1(x)=\begin{vmatrix}
        0 & e^{x}\sin(x)\\
        1 & e^{x}(\sin(x)+\cos(x))
    \end{vmatrix}=-e^{x}\sin(x)\] 
    \[W_2(x)=\begin{vmatrix}
        e^{x}\cos(x) & 0\\
        e^{x}(\cos(x)-\sin(x)) & 1
    \end{vmatrix}=e^{x}\cos(x)\]

    A partir da fórmula de variação dos parâmetros,
    \[y_p(x)=e^x\cos(x)\int\frac{\frac{e^x\sec(x)}{3}(-e^x\sin(x))}{e^{2x}}dx+e^x\sin(x)\int\frac{\frac{e^x\sec(x)}{3}(e^x\cos(x))}{e^{2x}}dx\]
    \[y_p(x)=-\frac{e^x\cos(x)}{3}\int\frac{\sin(x)}{\cos(x)}dx+\frac{e^x\sin(x)}{3}\int dx\]
    \[y_p(x)=\frac{e^x}{3}(\cos(x)\ln|\cos(x)|+x\sin(x))\]

    Ou seja, a solução geral da equação diferencial é
    \[\boxed{y(x)=c_1e^x\cos(x)+c_2e^x\sin(x)+\frac{e^x}{3}(\cos(x)\ln|\cos(x)|+x\sin(x))}\]

    \itemresolucao Resolvendo a equação homogênea associada,
    \[\lambda^2+2\lambda-8=0\Rightarrow(\lambda-2)(\lambda+4)=0\Rightarrow\lambda=2\text{ ou }\lambda=-4\]

    Nesse sentido, são soluções linearmente independentes da homogênea associada $y_1=e^{2x}$ e $y_2=e^{-4x}$, com $y_h=c_1e^{2x}+c_2e^{-4x}$. Calculando o Wronskiano dessas soluções,
    \[W(x)=\begin{vmatrix}
        e^{2x} & e^{-4x}\\
        2e^{2x} & -4e^{-4x}
    \end{vmatrix}=-6e^{-2x}\]

    Além disso,
    \[W_1(x)=\begin{vmatrix}
        0 & e^{-4x}\\
        1 & -4e^{-4x}
    \end{vmatrix}=-e^{-4x}\quad e \quad W_2(x)=\begin{vmatrix}
        e^{2x}  & 0\\
        2e^{2x} & 1
    \end{vmatrix}=e^{2x}\]

    Da fórmula de variação dos parâmetros,
    \[y_p(x)=e^{2x}\int\frac{(2e^{-2x}-e^{-x})(-e^{-4x})}{-6e^{-2x}}dx+e^{-4x}\int\frac{(2e^{-2x}-e^{-x})(e^{2x})}{-6e^{-2x}}dx\]
    \[y_p(x)=\frac{1}{9}e^{-x}-\frac{1}{4}e^{-2x}\]

    Assim, a solução geral da equação diferencial é
    \[\boxed{c_1e^{2x}+c_2e^{-4x}+\frac{1}{9}e^{-x}-\frac{1}{4}e^{-2x}}\]

    Agora, a fim de resolver o PVI, devemos substituir as condições iniciaisi fornecidas:
    \[y(0)=1\Rightarrow c_1+c_2-\frac{5}{36}=1\Rightarrow c_1+c_2=\frac{41}{36}\]
    \[y'(0)=0\Rightarrow 2c_1-4c_2-\frac{1}{9}+\frac{1}{2}=0\Rightarrow c_1-2c_2=-\frac{7}{36}\]

    Dessa forma, temos um sistema em $c_1$ e $c_2$:
    \[\begin{cases}
        c_1+c_2=\frac{41}{36}\\
        c_1-2c_2=-\frac{7}{36}
    \end{cases}\]
    Resolvendo esse sistema, encontramos $c_1=\frac{25}{36}$ e $c_2=\frac{4}{9}$. Ou seja, a solução geral do PVI é dada por
    \[\boxed{\frac{25}{36}e^{2x}+\frac{4}{9}e^{-4x}+\frac{1}{9}e^{-x}-\frac{1}{4}e^{-2x}}\]

\end{resolucao}

\begin{resolucao}
    Em primeira análise, percebemos que a derivada de maior ordem é multiplicada por um termo diferente de $1$, isto é, não está na forma normal. A fim de prosseguir com a resolução do problema, reorganizamos a equação nessa forma:
    \[y''+\frac{1}{x}y'+\left(1-\frac{1}{4x^2}\right)y=x^{-\frac{1}{2}}\]

    O enunciado indica que as funções dadas são soluções L.I. da equação homogênea associada. Nesse sentido, temos que a combinação linear dessas soluções também é solução da equação homogênea, da forma $y_h(x)=c_1x^{-\frac{1}{2}}\cos(x)+c_2x^{-\frac{1}{2}}\sin(x)$.

    Além disso, fornecidas as soluções, podemos encontrar seu Wronskiano:
    \[W(x)=\left|\begin{array}{c:c}
       x^{-\frac{1}{2}}\cos(x) & x^{-\frac{1}{2}}\sin(x)\\ \hdashline
       -\frac{1}{2}x^{-\frac{3}{2}}\cos(x)-x^{-\frac{1}{2}}\sin(x) & -\frac{1}{2}x^{-\frac{3}{2}}\sin(x)+x^{-\frac{1}{2}}\cos(x)
    \end{array}\right|=x^{-1}\]

    Além disso,
    \[W_1(x)=\left|\begin{array}{c:c}
        0 & x^{-\frac{1}{2}}\sin(x)\\ \hdashline
        1 & -\frac{1}{2}x^{-\frac{3}{2}}\sin(x)+x^{-\frac{1}{2}}\cos(x)
    \end{array}\right|=-x^{-\frac{1}{2}}\sin(x)\]
    \[W_2(x)=\left|\begin{array}{c:c}
        x^{-\frac{1}{2}}\cos(x)  & 0\\ \hdashline
        -\frac{1}{2}x^{-\frac{3}{2}}\cos(x)-x^{-\frac{1}{2}}\sin(x) & 1
    \end{array}\right|=x^{-\frac{1}{2}}\cos(x)\]

    Aplicando a fórmula de variação dos parâmetros, com $f(x)=x^{-\frac{1}{2}}$,
    \[y_p(x)=x^{-\frac{1}{2}}\cos(x)\int\frac{x^{-\frac{1}{2}}(-x^{-\frac{1}{2}}\sin(x))}{x^{-1}}dx+x^{-\frac{1}{2}}\sin(x)\int\frac{x^{-\frac{1}{2}}\cdot x^{-\frac{1}{2}}\cos(x)}{x^{-1}}dx\]
    \[y_p(x)=x^{-\frac{1}{2}}\cos^2(x)+x^{-\frac{1}{2}}\sin^(x)=x^{-\frac{1}{2}}\]

    Dessa maneira, a solução geral da equação diferencial é
    \[\boxed{y(x)=c_1x^{-\frac{1}{2}}\cos(x)+c_2x^{-\frac{1}{2}}\sin(x)+x^{-\frac{1}{2}}}\]

\end{resolucao}

\begin{resolucao}
    
    \itemresolucao Primeiramente, determinamos as soluções da equação homogênea associada:
    \[\lambda^3+\lambda=0\Rightarrow\lambda(\lambda^2+1)=0\Rightarrow\lambda=0,~\lambda=+i\text{ ou }\lambda=-i\]
    Dessa maneira, temos $y_1=1$, $y_2=\cos(x)$ e $y_3=\sin(x)$, com $y_h=c_1+c_2\cos(x)+c_3\sin(x)$. Assim, partimos para a determinação do Wronskiano das soluções:
    \[W(x)=\begin{vmatrix}
        1 & \cos(x) & \sin(x)\\
        0 & -\sin(x) & \cos(x)\\
        0 & -\cos(x) & -\sin(x)
    \end{vmatrix}=1\]
    \[W_1(x)=\begin{vmatrix}
        0 & \cos(x) & \sin(x)\\
        0 & -\sin(x) & \cos(x)\\
        1 & -\cos(x) & -\sin(x)
    \end{vmatrix}=1, 
    \quad W_2(x)=\begin{vmatrix}
        1 & 0 & \sin(x)\\
        0 & 0 & \cos(x)\\
        0 & 1 & -\sin(x)
    \end{vmatrix}=-\cos(x)\]
    \[W_3(x)=\begin{vmatrix}
        1 & \cos(x) &  0\\
        0 & -\sin(x) & 0\\
        0 & -\cos(x) & 1
    \end{vmatrix}=-\sin(x)\]

    A partir da fórmula de variação dos parâmetros, então,
    \[y_p(x)=1\int\tan(x)dx+\cos(x)\int\tan(x)(\cos(x))dx+\sin(x)\int\tan(x)(-\sin(x))dx\]
    \[y_p(x)=-\ln|\cos(x)|+\cos(x)\int\sin(x)dx+\sin(x)\int\frac{\sin^2(x)}{\cos(x)}dx\]

    Mas $\displaystyle\frac{\sin^2(x)}{\cos(x)}=\frac{1-\cos^2(x)}{\cos(x)}=\sec(x)-\cos(x)$. Dessa forma, temos

    \[y_p(x)=-\ln|\cos(x)|-\cos^2(x)+\sin(x)\left[\int\sec(x)dx-\int\cos(x)dx\right]\]
    \[y_p(x)=-\ln|\cos(x)|-\cos^2(x)+\sin(x)\left[\ln|\sec(x)+\tan(x)|-\sin(x)\right]\]
    \[y_p(x)=-\ln|\cos(x)|+\sin(x)\ln|\sec(x)+\tan(x)|-1\]

    Ou seja, a solução geral é dada por
    \begin{gather*}
        \boxed{y(x)=c_1'+c_2\cos(x)+c_3\sin(x)-\ln|\cos(x)|+\sin(x)\ln|\sec(x)+\tan(x)|},\\
        c_1'=c_1-1
    \end{gather*}

    \itemresolucao A equação característica da homogênea associada é 
    \[\lambda^3-2\lambda^2-\lambda+2=0\]
    
    Por teste, percebemos que $1$ é raiz da equação. Aplicando Briot-Ruffini para dividir a equação por $(\lambda-1)$,

    \begin{table}[H]
        \centering
        \begin{tabular}{c|cccc}
            1 & 1 & -2 & -1 & 2\\
            \hline
              & 1 & -1 & -2 & 0
        \end{tabular}
    \end{table}

    Dessa forma podemos fatorar a equação em $(\lambda-1)(\lambda^2-\lambda-2)=0$, o que implica que $(\lambda-1)(\lambda-2)(\lambda+1)=0$. Ou seja, são raízes da equação $\lambda=1$, $\lambda=2$ e $\lambda=-1$. Com isso, determinamos a solução homogênea como $y_h(x)=c_1e^x+c_2e^{2x}+c_3e^{-x}$. Calculando o Wronskiano das soluções,
    \[W(x)=\begin{vmatrix}
        e^x & e^{2x} & e^{-x}\\
        e^x & 2e^{2x} & -e^{-x}\\
        e^x & 4e^{2x} & e^{-x}
    \end{vmatrix}=6e^{2x}\]
    \[W_1(x)=\begin{vmatrix}
        0 & e^{2x} & e^{-x}\\
        0 & 2e^{2x} & -e^{-x}\\
        1 & 4e^{2x} & e^{-x}
    \end{vmatrix}=-3e^x, 
    \quad W_2(x)=\begin{vmatrix}
        e^x & 0 & e^{-x}\\
        e^x & 0 & -e^{-x}\\
        e^x & 1 & e^{-x}
    \end{vmatrix}=2\]
    \[W_3(x)=\begin{vmatrix}
        e^x & e^{2x} &  0\\
        e^x & 2e^{2x} & 0\\
        e^x & 4e^{2x} & 1
    \end{vmatrix}=e^{3x}\]

    Da fórmula de variação dos parâmetros, então,
    \[y_p(x)=e^x\int\frac{e^{4x}(-3e^x)}{6e^{2x}}dx+e^{2x}\int\frac{e^{4x}\cdot2}{6e^{2x}}dx+e^{-x}\int\frac{e^{4x}(e^{3x})}{6e^{2x}}dx\]
    \[y_p(x)=-\frac{e^x}{2}\int e^{3x}dx+\frac{e^{2x}}{3}\int e^{2x}dx+\frac{e^{-x}}{6}\int e^{5x}dx\]
    \[y_p(x)=\frac{1}{30}e^{4x}\]

    Assim, combinando as soluções homogênea e particular, temos
    \[\boxed{y(x)=c_1e^x+c_2e^{2x}+c_3e^{-x}+\frac{1}{30}e^{4x}}\]

\end{resolucao}

\begin{resolucao}
    Resolvendo a equação homogênea associada,
    \[\lambda^2+7\lambda=0\Rightarrow\lambda=\pm i\sqrt{7}\]
    Assim, tomamos $y_1=\cos(\sqrt{7}t)$, $y_2=\sin(\sqrt{7}t)$ e $y_h(t)=c_1\cos(\sqrt{7}t)+c_2\sin(\sqrt{7}t)$.

    Calculando o Wronskiano das soluções encontradas,
    \[W(t)=\begin{vmatrix}
        \cos(\sqrt{7}t)          & \sin(\sqrt{7}t)\\
        -\sqrt{7}\sin(\sqrt{7}t) & \sqrt{7}\cos(\sqrt{7}t)
    \end{vmatrix}=\sqrt{7}\]

    Além disso,
    \[W_1(t)=\begin{vmatrix}
        0 & \sin(\sqrt{7}t)\\
        1 & \sqrt{7}\cos(\sqrt{7}t)
    \end{vmatrix}=-\sin(\sqrt{7}t)\]
    \[W_2(t)=\begin{vmatrix}
        \cos(\sqrt{7}t)          & 0\\
        -\sqrt{7}\sin(\sqrt{7}t) & 1
    \end{vmatrix}=\cos(\sqrt{7}t)\]

    Da fórmula de variação de parâmetros, então,
    \[ y_p(t)=\cos(\sqrt{7}t)\int\frac{\sum_{k=3}^{n}e^{3kt}(-\sin(\sqrt{7}t))}{\sqrt{7}}dt+\sin(\sqrt{7}t)\int\frac{\sum_{k=3}^{n}e^{3kt}\cos(\sqrt{7}t)}{\sqrt{7}}dt\]

    Ou, reorganizando a fim de facilitar a integração,
    \[ y_p(t)=-\frac{\cos(\sqrt{7}t)}{\sqrt{7}}\sum_{k=3}^{n}\color{blue}\underbrace{\int e^{3kt}\sin(\sqrt{7}t)dt}_{\circled{1}}\color{black}+\frac{\sin(\sqrt{7}t)}{\sqrt{7}}\sum_{k=3}^{n}\color{red}{\underbrace{\int e^{3kt}\cos(\sqrt{7}t)dt}_{\circled{2}}}\]

    Em um primeiro momento, calcularemos a integral \circled{1} por partes. Tomamos $\sin(\sqrt{7}t)=u\rightarrow\sqrt{7}\cos(\sqrt{7}t)dt=du$ e $e^{3kt}dt=dv\rightarrow\dfrac{e^{3kt}}{3k}=v$. Dessa forma,
    \[\int e^{3kt}\sin(\sqrt{7}t)dt=\frac{e^{3kt}}{3k}\sin(\sqrt{7}t)-\frac{\sqrt{7}}{3k}\int e^{3kt}\cos(\sqrt{7}t)dt\]

    Novamente, calcularemos a integral resultante por partes. Tomamos $\cos(\sqrt{7}t)=u\rightarrow-\sqrt{7}\sin(\sqrt{7}t)dt=du$ e $e^{3kt}dt=dv\rightarrow\dfrac{e^{3kt}}{3k}=v$. Ou seja,
    \[\int e^{3kt}\sin(\sqrt{7}t)dt=\frac{e^{3kt}}{3k}\sin(\sqrt{7}t)-\frac{\sqrt{7}}{3k}\left[\frac{e^{3kt}}{3k}\cos(\sqrt{7}t)+\frac{\sqrt{7}}{3k}\int e^{3kt}\sin(\sqrt{7}t)dt\right]\]
    \[\int e^{3kt}\sin(\sqrt{7}t)dt=\frac{e^{3kt}}{3k}\sin(\sqrt{7}t)-\frac{\sqrt{7}e^{3kt}}{9k^2}\cos(\sqrt{7}t)-\frac{7}{9k^2}\int e^{3kt}\sin(\sqrt{7}t)dt\]
    \[\left(1+\frac{7}{9k^2}\right)\int e^{3kt}\sin(\sqrt{7}t)dt=\frac{e^{3kt}}{3k}\sin(\sqrt{7}t)-\frac{\sqrt{7}e^{3kt}}{9k^2}\cos(\sqrt{7}t)\]
    \[\color{blue}\text{\circled{1} }\int e^{3kt}\sin(\sqrt{7}t)dt\color{black}=\frac{e^{3kt}}{7+9k^2}\left[3k\sin(\sqrt{7}t)-\sqrt{7}\cos(\sqrt{7}t)\right]\]

    É feito procedimento análogo para a integral \circled{2}, utilizando integração por partes:
    \[\int e^{3kt}\cos(\sqrt{7}t)dt=\frac{e^{3kt}}{3k}\cos(\sqrt{7}t)+\frac{\sqrt{7}}{3k}\int e^{3kt}\sin(\sqrt{7}t)dt\]
    \[\int e^{3kt}\cos(\sqrt{7}t)dt=\frac{e^{3kt}}{3k}\cos(\sqrt{7}t)+\frac{\sqrt{7}}{3k}\left[\frac{e^{3kt}}{3k}\sin(\sqrt{7}t)-\frac{\sqrt{7}}{3k}\int e^{3kt}\cos(\sqrt{7}t)dt\right]\]
    \[\int e^{3kt}\cos(\sqrt{7}t)dt=\frac{e^{3kt}}{3k}\cos(\sqrt{7}t)+\frac{\sqrt{7}e^{3kt}}{9k^2}\sin(\sqrt{7}t)-\frac{7}{9k^2}\int e^{3kt}\cos(\sqrt{7}t)dt\]
    \[\left(1+\frac{7}{9k^2}\right)\int e^{3kt}\cos(\sqrt{7}t)dt=\frac{e^{3kt}}{3k}\cos(\sqrt{7}t)+\frac{\sqrt{7}e^{3kt}}{9k^2}\sin(\sqrt{7}t)\]
    \[\color{red}\text{\circled{2} }\int e^{3kt}\cos(\sqrt{7}t)dt\color{black}=\frac{e^{3kt}}{7+9k^2}\left[3k\cos(\sqrt{7}t)+\sqrt{7}\sin(\sqrt{7}t)\right]\]

    Substituindo as expressões encontradas na equação de $y_p$, 
    \begin{gather*}
        y_p(t)=-\frac{\cos(\sqrt{7}t)}{\sqrt{7}}\sum_{k=3}^{n}\frac{e^{3kt}}{7+9k^2}\left[3k\sin(\sqrt{7}t)-\sqrt{7}\cos(\sqrt{7}t)\right]+\\+\frac{\sin(\sqrt{7}t)}{\sqrt{7}}\sum_{k=3}^{n}\frac{e^{3kt}}{7+9k^2}\left[3k\cos(\sqrt{7}t)+\sqrt{7}\sin(\sqrt{7}t)\right]
    \end{gather*}
    \begin{gather*}
        y_p(t)=\sum_{k=3}^{n}\frac{e^{3kt}}{7+9k^2}\left[-\frac{3k}{\sqrt{7}}\cos(\sqrt{7}t)\sin(\sqrt{7}t)+\cos^2(\sqrt{7}t)\right.+\\\left.+\frac{3k}{\sqrt{7}}\sin(\sqrt{7}t)\cos(\sqrt{7}t)+\sin^2(\sqrt{7}t)\right]
    \end{gather*}
    \[y_p(t)=\sum_{k=3}^{n}\frac{e^{3kt}}{7+9k^2}\]

    Assim, encontramos a solução geral da EDO, da forma $y(t)=y_h(t)+y_p(t)$:
    \[\boxed{y(t)=c_1\cos(\sqrt{7}t)+c_2\sin(\sqrt{7}t)+\sum_{k=3}^{n}\frac{e^{3kt}}{7+9k^2}}\]

\end{resolucao}

\begin{resolucao}
    Considerando que o circuito é configurado em série, as fontes e quedas de tensões devem ser somadas. Nesse sentido, temos
    \[E(t)-\frac{1}{C}Q(t)-R\frac{dQ(t)}{dt}-L\frac{d^2Q(t)}{dt^2}=0\]
    \[500=\frac{1}{\frac{1}{108}\times10^{-3}}Q(t)+300\frac{dQ(t)}{dt}+0.2\frac{d^2Q(t)}{dt^2}\]
    \[Q''+1500Q'+540\cdot10^3Q=2500\]

    Resolvendo a equação homogênea associada,
    \[\lambda^2+1500\lambda+540\cdot10^3=0\]
    \[\lambda_{1,2}=\frac{1500\pm\sqrt{1500^2-4\cdot1\cdot540\cdot10^3}}{2}\Rightarrow\lambda_1=-600,~\lambda_2=-900\]

    Sendo assim, são soluções do sistema $Q_1(t)=e^{-600t}$ e $Q_2(t)=e^{-900t}$, e a solução homogênea pode ser expressa por uma combinação linear das duas, da forma $Q_h(t)=c_1e^{-600t}+c_2e^{-900t}$.

    Encontradas as soluções que resolvem a equação homogênea, prosseguimos com o cálculo do Wronskiano:
    \[W(t)=\begin{vmatrix}
        e^{-600t}          & e^{-900t}\\
        -600e^{-600t} & -900e^{-900t}
    \end{vmatrix}=-300e^{-1500t}\]

    Além disso,
    \[W_1(t)=\begin{vmatrix}
        0 & e^{-900t}\\
        1 & -900e^{-900t}
    \end{vmatrix}=-e^{-900t}\quad e \quad W_2(t)=\begin{vmatrix}
        e^{-600t} & 0\\
        -600e^{-600t} & 1
    \end{vmatrix}=e^{-600t}\]

    Aplicando a fórmula de variação dos parâmetros, 
    \[Q_p(t)=e^{-600t}\int\frac{2500(-e^{-900t})}{-300e^{-1500t}}dt+e^{-900t}\int\frac{2500e^{-600t}}{-300e^{-1500t}}dt\]
    \[Q_p(t)=\frac{25}{3}e^{-600t}\int e^{600t}dt-\frac{25}{3}e^{-900t}\int e^{900t}dt\]
    \[Q_p(t)=\frac{25}{1800}-\frac{25}{2700}=\frac{1}{216}\]

    Dessa forma, dado que $Q(t)=Q_h(t)+Q_p(t)$, encontramos que
    \[Q(t)=c_1e^{-600t}+c_2e^{-900t}+\frac{1}{216}\]

    O enunciado fornece as seguintes condições iniciais:
    \[\begin{cases}
        \displaystyle Q(0)=\frac{217}{216}\\
        \displaystyle i(0)=\frac{dQ}{dt}(0)=0
    \end{cases}\]

    Aplicando essas condições à expressão de $Q(t)$ encontrada,
    \[\begin{cases}
        \displaystyle\frac{217}{216}=c_1+c_2+\frac{1}{216}\\
        \displaystyle0=-600c_1-900c_2
    \end{cases}\Rightarrow
    \begin{cases}
        \displaystyle1=c_1+c_2\\
        \displaystyle0=c_1+\frac{3}{2}c_2
    \end{cases}\]

    Resolvendo o sistema, encontramos $c_1=3$ e $c_2=-2$. Dessa maneira, temos como resposta do problema
    \[\boxed{Q(t)=3e^{-600t}-2e^{-900t}+\frac{1}{216}}\]
    \[\boxed{i(t)=\frac{dQ}{dt}(t)=1800\left(e^{-900t}-e^{-600t}\right)}\]
\end{resolucao}

\end{document}